\subsection{মাধ্যম নির্বাচন: কারণসহ সিদ্ধান্ত}
\begin{longtable}{|p{\dimexpr.27\linewidth-7pt\relax}|p{\dimexpr.28\linewidth-7pt\relax}|p{\dimexpr.45\linewidth-7pt\relax}|}
\hline
\textbf{প্রয়োজন} & \textbf{সম্ভাব্য নির্বাচন} & \textbf{যুক্তি ও শর্ত}\\\hline
\endfirsthead
\hline
\textbf{প্রয়োজন} & \textbf{সম্ভাব্য নির্বাচন} & \textbf{যুক্তি ও শর্ত}\\\hline
\endhead
এক কক্ষের ল্যাব & UTP + Switch & কম খরচে সহজ সংযোগ; Ethernet চ্যানেলের অনুমোদিত দূরত্ব ও কেবলের শ্রেণি মানতে হবে।\\\hline
কারখানা বা দূরের ভবন & Optical Fiber & বেশি বহনক্ষমতা, কম ক্ষয় ও EMI-মুক্ত আলোকপথ; জোড়া লাগানো এবং উপযুক্ত ট্রান্সসিভার দরকার।\\\hline
কক্ষে চলমান ল্যাপটপ & Wi-Fi & কেবল ছাড়াই চলাচল; AP-এর আওতা, ব্যবহারকারীর সংখ্যা ও নিরাপত্তা বিবেচ্য।\\\hline
দুই দূরবর্তী অফিস & Fiber / পরিকল্পিত Radio Link & দূরত্ব, বাধা, অনুমোদিত ফ্রিকোয়েন্সি, টাওয়ার ও খরচ বিবেচ্য; শুধু প্রযুক্তির নাম যথেষ্ট নয়।\\\hline
স্যাটেলাইটভিত্তিক দূরাঞ্চল & VSAT / Satellite Terminal & স্থলসংযোগ কঠিন হলে উপযোগী; বিলম্ব, বৃষ্টি ও খরচের সীমাবদ্ধতা থাকে।\\\hline
TV Remote & Infrared & কাছের যন্ত্রে নির্দেশ পাঠায়; বাধা ও আলোকপথ বিবেচ্য।\\\hline
Cable TV & Coaxial Cable & সুরক্ষা আবরণযুক্ত পথে ভিডিও সংকেত বিতরণ; কানেক্টর ও কেবলের মান মিলতে হয়।\\\hline
\end{longtable}
\textbf{একই নেটওয়ার্কে একাধিক মাধ্যমের প্রয়োজন:} কলেজের কম্পিউটার থেকে সুইচ পর্যন্ত UTP, ভবনের মধ্যে ফাইবার এবং শিক্ষার্থীদের ল্যাপটপে Wi-Fi ব্যবহার করা যায়। প্রতিটি অংশের দূরত্ব, বহনক্ষমতা ও চলাচলের চাহিদা পৃথকভাবে মেটে। কেবল ব্যবহারে অনুমতিহীন প্রবেশের ঝুঁকি কমতে পারে। তবুও পরিচয় যাচাই ও তথ্যের গোপন সংকেতায়ন বা এনক্রিপশন প্রয়োজন হতে পারে।


\subsection{মাধ্যমের সংখ্যাগত তথ্য: দূরত্ব ও সময় আলাদা}
\chTwoPair{\textbf{একক ও Signal-এর গতি}\par
১ km = ১০০০ m; ১ $\mu$s = $10^{-6}$ s। শূন্যস্থানে বিদ্যুৎচৌম্বক তরঙ্গের গতি প্রায় $3\times10^8$ m/s। ফাইবারে আলোর গতি সাধারণ উদাহরণে প্রায় $2\times10^8$ m/s ধরা যায়; উপাদানভেদে বদলায়।}{\textbf{একটি পথের সময়}\par
প্রশ্নে Fiber-এর গতি $2\times10^8$ m/s এবং দূরত্ব ২ km দিলে শুধু Propagation Time $=2000/(2\times10^8)=10\,\mu$s। এটি শুধু সংকেত ছড়িয়ে পৌঁছানোর সময়; পুরো ফাইল পাঠানোর সময় নয়। মোট সময়ে ফাইলের বিটসংখ্যা ও ডেটা হার অনুযায়ী প্রেরণসময়, যন্ত্রে অপেক্ষা এবং প্রক্রিয়াকরণের সময়ও যোগ হয়।}
UTP চ্যানেল, কোঅ্যাক্সিয়াল অংশ ও ফাইবার মডিউলের কার্যকর দূরত্ব তারযুক্ত মাধ্যমের অংশে আছে। Bluetooth, Wi-Fi ও WiMAX-এর পরিসর বেতার যোগাযোগের অংশে দেওয়া হয়েছে। একটি সংযোগের দৈর্ঘ্য এবং পুরো LAN/WAN-এর বিস্তৃতি আলাদা। সংখ্যাগত তথ্য ব্যবহারের সময় এই পার্থক্য মানতে হয়।
